\chapter{Mixed Exercises}
\section{Worked Examples}
\noindent\textbf{Key Ideas and Methods.}\enspace
A mixed problem may become simpler after changing what is counted or optimized. The comparison count in randomized quicksort is a sum of indicators for pairs of keys being compared; these indicators need not be independent for linearity of expectation. Convexity and one-sided derivatives characterize minimizers of asymmetric losses, including distributions with atoms. For optimization involving infinite sums, establish a finite truncation bound before justifying the limiting step.

After finding a candidate, prove its feasibility, show that it attains the bound, and identify the equality conditions. An average-performance bound does not guarantee the same worst-case behavior, and existence of an optimum does not imply uniqueness.
\par\smallskip


\problem{CX01}{Comparisons in Randomized Quicksort}{}
Randomized quicksort sorts $n$ distinct keys. In each nonempty subproblem, it independently chooses a uniform pivot, compares it once with every other key in that subproblem, then recurses on the two sides. Let $C_n$ count all comparisons. Find $\E C_n$ exactly and explain why expected $O(n\log n)$ work does not mean that every run has this order of cost.
\solution Number the keys by increasing rank, $1,\ldots,n$. Let $I_{ij}$ indicate whether ranks $i<j$ are ever compared. Every comparison involves a pivot, which is then removed from future subproblems. A pair can therefore be compared at most once, and
\[
C_n=\sum_{1\le i<j\le n}I_{ij}.
\]
Consider the first pivot selected from ranks $\{i,i+1,\ldots,j\}$. Pivots outside this interval cannot separate $i$ from $j$. If the first pivot inside is an interior rank, it separates the two permanently. If it is either endpoint, the endpoints are compared then. Hence
\[
\Prob(I_{ij}=1)=\frac2{j-i+1}.
\]
Uniformity of the first pivot can be seen through an equivalent implementation: give all keys independent continuous random priorities and always choose the lowest-priority key in a subproblem. This has the same pivot law, and each rank in the interval is equally likely to have its lowest priority.

Linearity of expectation does not require the indicators to be independent. Grouping by $d=j-i$ and writing $H_n=\sum_{k=1}^n1/k$ gives
\[
\begin{aligned}
\E C_n
&=2\sum_{d=1}^{n-1}\frac{n-d}{d+1}\\
&=2(n+1)(H_n-1)-2(n-1)\\
&=\boxed{2(n+1)H_n-4n}.
\end{aligned}
\]
For $n=1$ the formula yields zero; define $C_0=0$ separately. For example, $\E C_4=29/6$.

If every pivot happens to be the smallest key in its subproblem, the cost is $(n-1)+\cdots+1=n(n-1)/2$. This event has positive probability $1/n!$. The expected order thus averages over random runs and is not a worst-case bound for each run.
\insight Ask when a particular pair of keys meets, rather than tracking the whole random recursion tree. Correlated pair indicators can still be added in expectation.

\problem{CX02}{Balancing Shortage and Excess Inventory}{}
Nonnegative demand $D$ has finite mean and distribution function $F$. A real order quantity $q\ge0$ incurs shortage cost $c_u>0$ per unit and excess cost $c_o>0$ per unit. Minimize
\[
L(q)=c_u\E(D-q)_++c_o\E(q-D)_+.
\]
Without assuming a continuous demand distribution, characterize every minimizer. For $\Prob(D=0,1,2)=(1/4,1/2,1/4)$, find all minimizers when $(c_u,c_o)=(3,1)$ and when $(c_u,c_o)=(2,1)$.
\solution For each fixed demand, the loss is a convex, piecewise linear function of $q$, so its expectation is convex. One-sided difference quotients are bounded by $\max(c_u,c_o)$, allowing their limits inside the expectation. Thus
\[
\begin{aligned}
L'_-(q)&=(c_u+c_o)F(q-)-c_u,\\
L'_+(q)&=(c_u+c_o)F(q)-c_u.
\end{aligned}
\]
At $q>0$, convexity makes nonpositive left derivative and nonnegative right derivative necessary and sufficient. Writing $\tau=c_u/(c_u+c_o)$, the full minimizer set is
\[
\boxed{\{q\ge0:F(q-)\le\tau\le F(q)\}}.
\]
At the boundary $q=0$, only the right derivative must be nonnegative. Since nonnegative demand implies $F(0-)=0\le\tau$, the same formula remains valid. Extending the loss formally to $q<0$ makes it strictly decreasing there, so the constraint loses no additional candidate.

A minimizer exists: $\tau\in(0,1)$, and the finite generalized quantile $\inf\{x:F(x)\ge\tau\}$ satisfies the displayed condition. If $F$ is flat at level $\tau$, an entire interval can minimize the loss; uniqueness must not be assumed.

For the discrete example and costs $(3,1)$, $\tau=3/4$. The distribution function jumps to $3/4$ at $1$ and remains there until $2$, yielding all minimizers $[1,2]$, each with expected loss $1$. For costs $(2,1)$, $\tau=2/3$ lies strictly inside the jump at $1$. The unique minimizer is then $q=1$, with minimum loss $3/4$.
\insight Asymmetric absolute loss selects a cost-weighted quantile. The mean is generally not optimal, and atoms or gaps in the distribution determine whether the optimal action is unique.

\problem{CX03}{Maximum Entropy with a Fixed Mean}{}
A nonnegative integer-valued random variable $N$ has mean $\mu>0$. With natural logarithms, define
\[
H(N)=-\sum_{k=0}^\infty p_k\log p_k,\qquad 0\log0=0.
\]
Prove a finite entropy bound and identify the unique maximizing distribution and its entropy. You may use the log-sum inequality for finite distributions, but must handle the infinite support without assuming finite entropy in advance.
\solution Put $r=\mu/(1+\mu)\in(0,1)$ and consider
\[
q_k=(1-r)r^k,\qquad k\ge0.
\]
This geometric distribution has mean $r/(1-r)=\mu$, and its entropy is
\[
H(q)=-\log(1-r)-\mu\log r
=(1+\mu)\log(1+\mu)-\mu\log\mu.
\]

For a rigorous general bound, retain states $0,\ldots,M-1$ separately and merge $N\ge M$ into one tail state of probability $t_M$. The reference distribution assigns probability $r^M$ to this tail. The finite log-sum inequality bounds the grouped entropy by
\[
H_M\le-\sum_{k<M}p_k\log q_k-t_M\log(r^M).
\]
The right-hand side is
\[
-(1-t_M)\log(1-r)
-\log r\left(\sum_{k<M}kp_k+Mt_M\right).
\]
As $M\to\infty$, $t_M\to0$ and the parenthesized quantity is $\E\min(N,M)\uparrow\mu$. Also
$H_M=-\sum_{k<M}p_k\log p_k-t_M\log t_M$ tends to $H(N)$, even if the latter is initially allowed to be infinite, since $-t\log t\to0$. Consequently
\[
H(N)\le(1+\mu)\log(1+\mu)-\mu\log\mu<\infty.
\]

To prove uniqueness, the difference between the grouped cross-entropy bound and $H_M$ is the grouped relative entropy $D_M\ge0$. Its limit is $H(q)-H(N)$. If $p\ne q$, select $j$ with $p_j\ne q_j$. For every $M>j$, merge all groups other than state $j$. The log-sum inequality gives
\[
D_M\ge
p_j\log\frac{p_j}{q_j}
+(1-p_j)\log\frac{1-p_j}{1-q_j}>0.
\]
This positive lower bound does not depend on $M$, excluding equality. Thus the geometric distribution uniquely maximizes entropy.
\insight On an infinite state space, first aggregate the tail and then pass to the limit. This avoids presupposing finite entropy when manipulating relative entropy.


\section{Exercises}
\noindent\textbf{Key Ideas and Methods.}\enspace
Translate each constraint into a structure that can be computed. Congruence conditions suggest residue classes or roots-of-unity filters. Rotational symmetry can reduce conditional moments to averages over a sphere, with attention to the sampling measure. Patterns in repeated trials can be tracked by finite states recording the relevant suffix. Variance calculations must retain covariance terms caused by shared variables or dependence over time.

Integer problems often require separating prime exponents from the remaining factors. Division in modular arithmetic is valid only when the divisor is invertible. Asymptotic counts must account for boundary terms and repeated entries. Finite enumeration can check small cases; a general conclusion still needs a recurrence, identity, or invariant that proves it.
\par\smallskip



\exercise{AA01}{Local Independence with a Global Parity Constraint}
Generate $n-1$ independent fair bits, $n\ge3$, and set $B_n$ equal to their sum modulo 2. Prove that every set of $n-1$ bits is independent and fair, whereas all $n$ bits are not independent. Find the probability that they are all zero.

\exercise{AA02}{The Number of Successes Modulo Three}
Toss $n\ge0$ fair coins independently and let $S$ count heads. Find $\Prob(S\equiv r\pmod3)$ for $r=0,1,2$. You may use $\omega=e^{2\pi i/3}$, but justify the root-of-unity filter. Evaluate $n=8,r=0$.

\exercise{AA19}{Return Probability under Adjacent Swaps}
Three distinct cards start in order 123. At each step, independently swap positions 1 and 2 or positions 2 and 3 with equal probability. Find the probability of returning to 123 after $t\ge0$ steps and explain why it has no single long-time limit.

\exercise{AA22}{Isolated Marks in a Circular Arrangement}
Randomly arrange $m$ A's and $n$ B's on $N=m+n\ge3$ labeled circular positions, uniformly over type arrangements. A B is isolated if both neighbors are A. Find the expected isolated-B count, then the corresponding expectation on a line counting only nonendpoint centers. Evaluate $m=4,n=3$.

\exercise{AA03}{A Projection of a Uniform Disk Point}
Choose $(X,Y)$ uniformly by area in a disk of radius $R>0$. Find the density of $X$, $\E|X|$, and $\Var X$, and explain why $X$ is not uniform on $[-R,R]$.

\exercise{AA07}{Linear Separation of Two Classes on a Circle}
Choose $m+n$ independent uniform points on the unit circle. Label the first $m\ge2$ as A and the remaining $n\ge2$ as B. Find the probability that a line strictly separates all A points from all B points. Evaluate $m=n=3$.

\exercise{AA09}{A Gaussian Ratio in the Positive Quadrant}
Let $X,Y$ be independent standard normals. Conditional on $X>0,Y>0$, define $Z=X/Y$. Find its distribution, density, and median, and decide whether its conditional mean is finite.

\exercise{AA10}{Dependence of Squared Coordinates at Fixed Energy}
Let $X,Y,Z$ be independent standard normals and $S=X^2+Y^2+Z^2$. Find $\E[X^2Y^2\mid S]$, $\Var(X^2\mid S)$, and $\Cov(X^2,Y^2\mid S)$. Derive the fourth moments from rotation symmetry instead of citing a sphere-moment formula.

\exercise{AA13}{The Closest Gap among Three Arrival Times}
Three arrival times are independent and uniform on $[0,L]$, with $L>0$. Let $D$ be the smallest time difference between any two arrivals. Find its complete distribution function and mean; distances to the interval endpoints are not included.

\exercise{AA16}{The Minimum and Range of an Exponential Sample}
Let $X_1,\ldots,X_n$ be independent exponentials of rate $\lambda>0$, with $n\ge2$. Put $U=\min_iX_i$, $V=\max_iX_i$, and $D=V-U$. Prove that $U,D$ are independent, find the distribution of $D$, and compute $\E[V\mid U]$.

\exercise{AA17}{Correlation after Squaring Gaussian Variables}
Let $(X,Y)$ be jointly standard normal with correlation $\rho\in[-1,1]$. Find $\E[Y^2\mid X]$, $\Cov(X^2,Y^2)$, and $\Corr(X^2,Y^2)$. Can the squared variables' correlation identify the original sign?

\exercise{AA20}{Reciprocal Distance to a Uniform Sphere Point}
In three dimensions, let $X$ be uniform by area on the sphere of radius $R>0$. A fixed point $a$ is distance $r\ge0$ from the center, and $D=|X-a|$. Find $\E(1/D)$. For $r>0$, also find $\E(1/D^2)$ and analyze $r=R$.

\exercise{AA06}{The First Repeated Adjacent Outcome}
Independent trials succeed with probability $p\in(0,1)$; put $q=1-p$. Starting before any trials, stop when two successive outcomes first agree. Find the mean total trial count and the probability of stopping on two successes.

\exercise{AA08}{Variance of an Adjacent Pattern Count}
Toss a coin with head probability $p\in(0,1)$ independently $n\ge2$ times. Let $N$ count adjacent occurrences of HT at all starting positions. Find its mean and variance, and evaluate $n=10,p=1/2$.

\exercise{AA12}{Temporal Correlation along a Cycle Walk}
Consider a symmetric nearest-neighbor walk on a cycle of $n\ge3$ vertices, with $X_0$ uniform over the vertices. Put $Y_t=\cos(2\pi X_t/n)$ and $S_T=\sum_{t=0}^{T-1}Y_t$, $T\ge1$. Find $\Cov(Y_0,Y_k)$, $\Var S_T$, and $\lim_{T\to\infty}\Var S_T/T$.

\exercise{AA04}{Inverting the Trailing-Zero Count}
For a nonnegative integer $K$, find every nonnegative integer $n$ for which $n!$ has exactly $K$ trailing decimal zeros, with $0!=1$. Avoid computing factorials. Prove that a nonempty solution set consists of five consecutive integers, and solve the cases $K=24$ and $K=29$.

\exercise{AA05}{The Last Nonzero Digit of a Binomial Coefficient}
Design an algorithm for the last nonzero decimal digit of $\binom nk$, for integers $0\le k\le n$, without first constructing the full coefficient. Separate powers of 2 and 5 from the unit factors. Explain why direct division modulo 10 is invalid, and evaluate $\binom{20}{7}$.

\exercise{AA11}{Visibility between Random Lattice Points}
Choose two distinct points uniformly from $\{0,\ldots,M\}\times\{0,\ldots,N\}$, with $M,N\ge1$. They are visible to each other if their connecting segment has no other integer lattice point in its interior. Give a finite-sum formula for the probability and evaluate $M=N=2$.

\pagebreak
\exercise{AA14}{The Asymptotic Number of Integer-Sided Triangles}
Let $T(N)$ count positive integer triples $x\le y\le z$ with $x+y+z=N$ and $x+y>z$. Prove $T(N)=N^2/48+O(N)$ and list all solutions for $N=12$. Explain why dividing the ordered count by 6 is generally not exact.

\exercise{AA15}{Algebraic Powers Close to Integers}
Let $\alpha=(3+\sqrt5)/2$. Construct an integer recurrence giving the nearest integer $L_n$ to $\alpha^n$ for every $n\ge1$. Prove uniqueness and find the exact error and parity pattern. Find the smallest $n$ for which the error is below $10^{-6}$.

\exercise{AA21}{Fixing Both Digit Sum and Digit Product}
Find the smallest positive decimal integer with digit product 72 and digit sum 14. Then change the digit sum to 18 and find the new minimum. Explain why a digit 1 cannot always be deleted in this problem.

\exercise{AA18}{Transcendental Logarithms and an Irrational Ratio}
Prove that $\log2$ is transcendental and that $\log2/\log3$ is irrational. Explain why the ratio of two transcendental numbers need not be transcendental. Here $\log$ is natural logarithm. You may use the Hermite--Lindemann theorem: if $z$ is nonzero and algebraic, then $e^z$ is transcendental.
